\documentclass[11pt]{article}
\usepackage[margin=25mm]{geometry}
\usepackage{fontspec}
\setmainfont{Latin Modern Roman}
\usepackage{amsmath,amssymb,mathtools}
\usepackage{unicode-math}
\setmathfont{Latin Modern Math}
\usepackage{xurl,hyperref}
\hypersetup{colorlinks=true,urlcolor=blue}
\setcounter{secnumdepth}{0}
\setlength{\parindent}{0pt}
\setlength{\parskip}{5pt}
\setlength{\emergencystretch}{3em}
\newcommand{\sref}[2]{\hyperlink{source:#1:#2}{[S#2]}}
\newcommand{\cataloguescope}[1]{\paragraph{Recorded target.}#1}
\newcommand{\R}{\mathbb R}
\newcommand{\C}{\mathbb C}
\newcommand{\N}{\mathbb N}
\newcommand{\Z}{\mathbb Z}
\newcommand{\Q}{\mathbb Q}
\newcommand{\F}{\mathbb F}
\newcommand{\eps}{\varepsilon}
\begin{document}
\section{From Dislocation Mechanics to Continuum Plasticity}\label{statement}

\subsection{Definitions and mathematical statement}

\paragraph{Shared notation.}$\mathbb N=\{1,2,\ldots\}$; the standard number fields have their usual meanings. All additional conventions are stated locally.


A dislocation is a lattice defect described microscopically by positions or curves with Burgers vectors. A continuum plasticity model replaces the discrete configuration by fields describing plastic deformation and dissipation. A scale transition must state the small-scale parameter, limiting energy or evolution, and topology of convergence.

\paragraph{Catalogue wording retained for traceability.}
Develop a rigorous scale-transition theory connecting discrete dislocation mechanics to continuum models of crystal plasticity.

\paragraph{Scope and source qualification.}This is a programme for deriving models; the source gives no single universal energy or convergence theorem to substitute for it.
\subsection{Short English statement}
Derive continuum crystal-plasticity descriptions rigorously from discrete dislocation mechanics.
\subsection{Research attempt}
\paragraph{Model, date and investigated scope.}
OpenAI GPT-6 Astra (configured model \texttt{gpt-6-astra}), ultra reasoning; generated 11 October 2026. Calculated the logarithmic dislocation self-energy and identified competing renormalizations that a continuum derivation must distinguish.

\paragraph{Sources and literature check.}
Read the full supplied definition, catalogue wording, source qualifications and listed references for this question. No fresh status or novelty search was performed. The argument below proves its auxiliary statements directly and does not import an unverified external theorem.

\paragraph{A single-defect core calculation.} In a two-dimensional scalar antiplane model let the Burgers magnitude be \(b\ne0\), core radius \(\epsilon\), and outer radius \(R>\epsilon\). The locally defined displacement \(u=b\theta/(2\pi)\) has circulation \(b\) and single-valued strain \(\nabla u=b e_\theta/(2\pi r)\). For shear modulus \(\mu>0\),
\[
 E_\epsilon=\frac\mu2\int_{\epsilon<|x|<R}|\nabla u|^2dx
 =\frac{\mu b^2}{4\pi}\log(R/\epsilon).
\]
Thus even one fixed dislocation has a divergent elastic self-energy as the core shrinks. A limit based only on weak convergence of the defect measure \(b\delta_0\) cannot preserve unrenormalized finite energy.
\paragraph{What a scale transition must retain.} Dividing by \(|\log\epsilon|\) leaves the coefficient \(\mu b^2/(4\pi)\) for this isolated defect. Subtracting that divergence instead retains finite outer-scale terms. These are different limits. For several defects, the cross term in \(|\sum_i\nabla u_i|^2\) additionally depends on their relative positions and signs, whereas each core term depends on \(b_i^2\). Losing either datum before taking the limit can lose physical energy.
\paragraph{First gap.} This calculation establishes neither compactness nor an energy-limit theorem for many dislocations and includes no motion law or dissipation. The source programme needs specified core physics, number and spacing scaling, topology, and an evolution connecting discrete motion to continuum plasticity. The first missing step is a uniform coercive energy-and-dissipation estimate for such a specified many-defect scaling. The scalar antiplane example is explicitly an auxiliary test, not a universal crystal model.

\paragraph{Verification and adversarial audit.}
Checked Burgers circulation, the elastic prefactor and logarithmic integral; distinguished local multivalued displacement from its single-valued strain and energy from evolution. Model self-audit only; no independent human mathematical verification. Applied contribution prompt A before drafting and B afterwards. Checked Burgers circulation, the elastic prefactor and logarithmic integral; distinguished local multivalued displacement from its single-valued strain and energy from evolution. First remaining gap: Prove a uniform compactness and energy-dissipation estimate in a specified many-defect/core scaling appropriate to the source crystal-plasticity programme. The inherited source status is preserved; original-target completion remains zero.

Coordinator checked the logarithmic antiplane core-energy coefficient and distinction between subtraction and division normalizations. This additional check was model-assisted, not independent human verification.

\paragraph{Tools, process and contributions.}
Python 3 and local file tools used to read definitions and serialize manuscripts. Calculations are symbolic reasoning; no numerical experiment or theorem prover. Contribution prompt A was followed by prompt B. GPT-6 Astra with ultra reasoning drafted the research and performed its stated self-audit; any additional coordinator contribution is identified in the verification and audit. A GPT-6-based Codex coordinator assembled the manuscript and checked repository integration; its exact variant and reasoning configuration were not exposed. No human mathematical contribution, independent human proof check or formal proof certification is supplied. Complete derivations and the consulted source references are included here; private reasoning transcripts are not published.

\paragraph{Final status and first remaining gap.}
\textbf{UNRESOLVED}. Prove a uniform compactness and energy-dissipation estimate in a specified many-defect/core scaling appropriate to the source crystal-plasticity programme.

\section{Completion Estimate}
\noindent\textbf{COMPLETION: 0\%}
This is an elementary auxiliary calculation or an explicit test of a restricted strategy. No new portion of the full source target is established.

\subsection{Sources}
\begin{itemize}
\item[\hypertarget{source:statement:1}{S1}] ULAM.AI and the UnsolvedMath Contributors, UnsolvedMath, catalogue entry. Catalogue curation under CC BY 4.0; definitions and English exposition added locally.
\url{https://huggingface.co/datasets/ulamai/UnsolvedMath}.
\item[\hypertarget{source:statement:2}{S2}] Mechanics of Materials: Multiscale Design of Advanced Materials and Structures (2024), Oberwolfach Reports, DOI 10.4171/owr/2024/11, pp. 657--730. Original report and mathematical context. Incorporated report material remains under the source CC BY-SA 4.0 terms unless otherwise noted.
\url{https://doi.org/10.4171/owr/2024/11}.
\end{itemize}
\label{end:statement}
\end{document}
